\section{Realizability Theorem}\label{sec:requirements}
%==============================================================================

This section asks when a concrete host system can realize the unit-rate corner isolated earlier and identifies the two structural properties needed for that purpose.

A host system realizes the rate-$1$ regime when one location is authoritative and every secondary representation is maintained as a deterministic view. Two capabilities are required:
\begin{enumerate}
\item \textbf{Causal update propagation}: updates to the authoritative source must automatically propagate to derived locations.
\item \textbf{Provenance observability}: the system must expose enough structural information to verify which locations are authoritative and which are derived.
\end{enumerate}

\subsection{Confusability and Side Information}\label{sec:confusability}

To connect realizability to information available at verification time, fix a fact $F$ and consider the alternatives that remain compatible with the available observations. These alternatives form the same zero-error obstruction analyzed in Sections~\ref{sec:converse}--\ref{sec:equality}: if the host cannot separate the surviving ambiguity class, exact verification requires additional structural side information.

\begin{lemma}[Confusability Clique Bound]\label{lem:confusability-clique}
If the confusability graph for $F$ contains a clique of size $k$, then any zero-error side-information scheme distinguishing those $k$ alternatives requires at least $k$ distinct tags, equivalently at least $\log_2 k$ bits of side information.
\end{lemma}
\leanmeta{\LH{OBS5}}

\begin{proof}
Within a $k$-clique, every pair of alternatives is confusable under the base observation. A zero-error tag assignment must therefore separate all $k$ states. Distinct states cannot share a tag, so at least $k$ tags are required.
\end{proof}

If a host system does not expose enough structural information to separate the remaining alternatives, exact verification of the rate-$1$ regime is impossible.

\subsection{The Structural Timing Constraint}\label{sec:timing}

\begin{definition}[Structural Fact]\label{def:structural-fact-req}
A fact $F$ is \emph{structural} if every location encoding $F$ is fixed when the underlying object, declaration, or artifact is created. After that moment, the encoding can be read and compared, but not retroactively re-created without a new edit event.
\end{definition}

Structural facts are the natural exact-consistency setting because they expose a clear creation moment. Examples include schema declarations, registry membership rules, dependency metadata, and interface-level constraints.

\begin{theorem}[Timing Constraint for Structural Derivation]\label{thm:timing-forces}
For structural facts, derivation must occur no later than the moment at which the relevant encoding locations become fixed.
\end{theorem}
\leanmeta{\LH{REQ1}, \LH{TRI1}}

\begin{proof}
Let $t_{\mathrm{fix}}$ denote the time at which the structural encoding becomes fixed. A derivation performed strictly before $t_{\mathrm{fix}}$ cannot depend on the completed source value; a derivation performed strictly after $t_{\mathrm{fix}}$ cannot alter the already fixed structural representation. Therefore structural derivation must occur at $t_{\mathrm{fix}}$ or as part of the same creation/update event.
\end{proof}

\subsection{Requirement 1: Causal Update Propagation}\label{sec:causal-propagation}

\begin{definition}[Causal Update Propagation]\label{def:causal-propagation}
An encoding system has \emph{causal update propagation} if every update to a source location is followed by updates to all locations derived from that source with no reachable intermediate state in which source and derived locations disagree. Equivalently, propagation is part of the same admissible update event for the purposes of reachable-state analysis, rather than an asynchronous repair step with observable stale states.
\end{definition}

This is the host-system manifestation of deterministic dependence. Without it, a temporal gap appears between source modification and derived-view repair.

\begin{theorem}[Causal Propagation is Necessary for Unit Rate]\label{thm:causal-necessary}
Achieving independent rate $1$ for structural facts requires causal update propagation.
\end{theorem}
\leanmeta{\LH{REQ2}}

\begin{proof}
By Theorem~\ref{thm:timing-forces}, structural derivation must occur at the moment the source-side structural encoding is fixed. If causal propagation is absent, then some derived location remains stale until a later manual action. During that interval the source and derived locations disagree, so the architecture does not realize exact consistency. Hence a verifiable rate-$1$ realization for structural facts requires causal propagation.
\end{proof}

\subsection{Requirement 2: Provenance Observability}\label{sec:provenance-observability}

\begin{definition}[Provenance Observability]\label{def:provenance-observability}
An encoding system has \emph{provenance observability} if it supports queries that reveal which locations encode a fact, which of those locations are authoritative, and which are derived.
\end{definition}

Provenance observability is the structural side information needed to certify that the system is genuinely in the single-source regime rather than merely appearing coherent on a small sample of states.

\begin{theorem}[Provenance Observability is Necessary for Verifiable Unit Rate]\label{thm:provenance-necessary}
Verifying that independent rate $1$ holds requires provenance observability.
\end{theorem}
\leanmeta{\LH{REQ3}}

\begin{proof}
To verify independent rate $1$, one must enumerate the locations encoding the fact, identify the authoritative source, and confirm that every remaining location is derived from it. Without provenance queries, these checks cannot be completed from inside the host system. The architecture may be coherent on observed states, but the single-source claim is not verifiable.
\end{proof}

\subsection{Independence of the Two Requirements}\label{sec:independence}

The two requirements are logically separate.

\begin{theorem}[Requirements are Independent]\label{thm:independence}
\begin{enumerate}
\tightlist
\item An encoding system can have causal propagation without provenance observability.
\item An encoding system can have provenance observability without causal propagation.
\end{enumerate}
\end{theorem}
\leanmeta{\LHrng{IND}{1}{2}}

\begin{proof}
For (1), consider a system that automatically refreshes all derived views after each source update but exposes no query interface for its derivation graph. Coherence may hold operationally, but the single-source claim is not inspectable. For (2), consider a system that records complete provenance metadata yet requires users to trigger propagation manually after each change. The derivation structure is visible, but stale views remain reachable. Hence neither property implies the other.
\end{proof}

\subsection{The Realizability Theorem}\label{sec:realizability-theorem}

\begin{theorem}[Necessary and Sufficient Realizability Conditions]\label{thm:ssot-iff}
An encoding system can achieve verifiable independent rate $1$ for structural facts if and only if it provides both causal update propagation and provenance observability.
\end{theorem}
\leanmeta{\LH{SOT1}}

\begin{proof}
Necessity follows from Theorems~\ref{thm:causal-necessary} and~\ref{thm:provenance-necessary}. For sufficiency, assume both properties hold. Causal propagation ensures that every derived location is updated as part of the same structural event as the source; provenance observability then certifies that all secondary encodings are in fact derived from that source. The architecture therefore realizes a verifiable single-source encoding, i.e., independent rate $1$.
\end{proof}

The next section records representative readings of the theorem on familiar host systems.

%==============================================================================
